\documentclass[10pt,a4paper]{amsart}
\usepackage[margin=19mm]{geometry}
\usepackage{amsmath,amssymb,mathtools}
\usepackage[scr=rsfs,cal=euler]{mathalfa}
\usepackage{xfrac}
\usepackage[unicode,hidelinks,bookmarks=false]{hyperref}
\newcommand{\ie}{i.e.\ }
\newcommand{\cf}{cf.\ }
\newcommand{\resp}{respectively\ }
\newcommand{\Subsection}{\subsection}
\providecommand{\nobreakdash}{\nobreak\hbox{-}}
\setlength{\emergencystretch}{2em}
\multlinegap0pt
\usepackage{bm}
\usepackage{bbm}
\usepackage{dsfont}
\usepackage{xspace}
\usepackage{cancel}
\usepackage{mathtools}
\usepackage{amsthm}
\makeatletter
\@ifundefined{theorem}{\newtheorem{theorem}{Theorem}}{}
\makeatother
\newcommand{\R}{\mathcal{R}}
\newcommand{\Z}{\mathbb{Z}}
\newcommand{\id}{\mathrm{id}}
\title[Reversibility, balance and expansivity of non-uniform cellul]{Reversibility, balance and expansivity of non-uniform cellular automata}
\author{Katariina Paturi}
\date{Source: arXiv:2507.06896v1 (CC BY 4.0)}
\begin{document}
\maketitle

\noindent\emph{Source and licence.} Reversibility, balance and expansivity of non-uniform cellular automata. Version arXiv:2507.06896v1 (CC BY 4.0), distributed under the Creative Commons Attribution 4.0 International licence, as recorded in the arXiv licence metadata for this exact version and shown on the abstract page https://arxiv.org/abs/2507.06896v1. Full rights notice: licenses/arxiv-2507.06896-CC-BY-4.0.txt.\par\smallskip

\noindent\textit{Editorial scope.} Counted unit: the Theorem whose source label is \texttt{unif\_revers\_unif} in the article (source0.tex lines 257--259, source label \texttt{unif\_revers\_unif}), delivered together with the article's own complete original proof of it (source0.tex lines 261--296, one \texttt{proof} environment of 36 lines). No mathematics is written, repaired or re-derived: statement and proof are the article's own text line for line. Required context retained uncounted: none. Every definition, display and label of those lines is retained unchanged. Omitted from this package and not redistributed: 1--256, 260--260, 297--774. Nothing inside a retained excerpt is omitted or altered: every retained body is delivered exactly as the article's lines are written, so no omission receipt of any kind accompanies this package. Citations: the retained excerpts carry no citation command at all, so no bibliography entry of the article is transcribed and no reference is redistributed. Reference adaptation: none was needed and none was made; every reference inside the delivered unit resolves inside it, so no unresolved reference or citation is produced. Printed slips kept literal: no printed word, symbol, hypothesis or number of the counted statement or of its proof was altered. Layout and class: the article's own class is \texttt{llncs} with packages including fontenc, graphicx, amsmath, amssymb, mathtools, enumitem; the wrapper is the lane's pinned \texttt{amsart} editorial wrapper at 10pt on A4, which restates the theorem-like environments the retained text needs and adds the article's own macro definitions that the retained text needs (\texttt{\textbackslash R}, \texttt{\textbackslash Z}, \texttt{\textbackslash id}), with extra wrapper packages (bm, bbm, dsfont, xspace, cancel, mathtools). The delivered numbering is the unit's own. Assets: no retained span contains a figure environment, a graphics-inclusion command, a PDF-page inclusion, a table or a TikZ picture; no image file is copied into this package, the package declares no asset dependency and no image file is redistributed with it. Licence: version arXiv:2507.06896v1 of this article is distributed under the Creative Commons Attribution 4.0 International licence, as recorded in the arXiv licence metadata for this exact version and shown on the abstract page https://arxiv.org/abs/2507.06896v1. A full rights notice is delivered at licenses/arxiv-2507.06896-CC-BY-4.0.txt. Characters: every retained excerpt is pure ASCII, so the delivered engine, XeLaTeX with its default Latin Modern text font, reports no missing character.
\begin{theorem} \label{unif_revers_unif}
Let $\R_1$ be a finite rule set and $\theta \in \R_1^\Z$ uniformly recurrent. Let $\R_2$ be a rule set containing no pair of different identical rules, and $\phi \in \R_2^\Z$ a distribution such that $H_\phi = H_\theta^{-1}$. Then $\phi$ is uniformly recurrent.
\end{theorem}

\begin{proof} 
Let $r\in \Z_+$ be at least as large as the radius of any rule in $\R_1$ or $\R_2$. We show that for any $x,y \in \Z$, if $\theta_{|[x-r,x+r]} = \sigma^k (\theta_{|[y-r,y+r]})$, where $k=y-x$, then $\phi(x) = \phi(y)$. For all $c \in \Sigma^\Z$ we have $H_\theta (c)_{[x-r,x+r]} = H_{\sigma^k (\theta)} (c)_{[x-r,x+r]}$. Then applying the local rule $\phi(x)$, we have
\begin{align*}
H_\phi (H_\theta (c)) (x) = H_\phi (H_{\sigma^k (\theta)} (c)) (x).
\end{align*}
Since $H_\phi \circ H_\theta = \id = H_{\sigma^k (\phi)} \circ H_{\sigma^k (\theta)}$, we then have
\begin{align*}
H_{\sigma^k(\phi)} (H_{\sigma^k (\theta)} (c)) (x) = H_\phi (H_{\sigma^k (\theta)} (c)) (x).
\end{align*}
Because $H_{\sigma^k (\theta)}$ is surjective, for any $e \in \Sigma^\Z$ there is $c \in \Sigma^\Z$ such that $H_{\sigma^k (\theta)}(c) = e$. Therefore for any $e\in \Sigma^\Z$, 
\begin{align*}
H_{\sigma^k(\phi)} (e) (x) = H_\phi (e) (x).
\end{align*}
Since there are no identical rules in $\R_2$, $\sigma^k(\phi)(x) = \phi(x)$. Now because $\theta$ is uniformly recurrent, for any $m \in \Z_+$, and any segment $D \subseteq \Z$ of length $|D|=m+2r$, there is $n \in \Z_+$ such that  a copy of $\theta_{|D}$ appears in every length $n$ segment of $\theta$. Then by the previous, for every segment $D'$ of length $|D'|=m$ there is $n\in \Z_+$ such that a copy of $\phi_{|D'}$ appears in every segment of length $n$. Hence $\phi$ is uniformly recurrent.
\qed

%Let $x,y \in \Z$, $x\leq y$ and $r\in \Z_+$ at least as large as the radius of any rule in $\R_1$ or $\R_2$. Because $\theta$ is uniformly recurrent, there is $n \in \Z_+$ such that a disjoint copy of the pattern $\theta_{|[x-2r,y+2r]}$ appears in every segment of length $n$ in $\theta$. Let $k\in \Z$ be such that $\theta_{|[x-2r+k,y+2r+k]}$ is a disjoint copy of $\theta_{|[x-2r,y+2r]}$. We show that $\phi_{|[x+k,y+k]}$ is a copy of $\phi_{|[x,y]}$.



%Let then $e \in \Sigma^\Z$ be such that $e(z) = e(z+k)$ and for all $z' \in [z-r,z+r]$, $H_\theta (e)(z') = H_\theta(e)(z'+k) = c(z')$. Such a configuration exists, because $H_\theta$ is bijective and the neighbourhood $N_\theta([z-r,z+r])$ doesn't overlap with the neighbourhood  $N_\theta([z-r+k,z+r+k])$. But now
%\begin{align*}
%H_\phi(e)(z) = H_\phi(c)(z) \neq H_\phi(\sigma^{-k}(c))(z+k) = H_\phi (e)(z+k),
%\end{align*}
%meaning $H_\phi \circ H_\theta (e) \neq e$, which is a contradiction. 

%Let $z \in [x,y]$, $c\in \Sigma^\Z$ and $e = H_\phi (c)$. Because $\theta_{|[x-2r+k,y+2r+k]}$ is a copy of $\theta_{|[x-2r,y+2r]}$, then 
%\begin{align*}
%H_\theta (\sigma^{-k}(e))_{|[x-r+k,y+r+k]} = \sigma^{-k}(c)_{|[x-r+k,y+r+k]}
%\end{align*}
%Now because $H_\phi = H_\theta^{-1}$, 
%\begin{align*}
%H_\phi (\sigma^{-k}(c)) (z+k) =  \sigma^{-k}(e)(z+k) = e(z) = H_\phi (c)(z).
%\end{align*}
%Therefore there's no pattern that $\phi(z)$ and $\phi(z+k)$ map differently, and hence they are identical. Because $\R_2$ contained no pairs of different identical rules, then $\phi(z) = \phi(z+k)$. Then there is $n \in \Z_+$ such that a copy of pattern $\phi_{|[x,y]}$ appears in every segment of length $n$ in $\phi$, and $\phi$ is uniformly recurrent.
\end{proof}

\end{document}
